\section{Lecture 22: Sequence Models (RNN, LSTM, GRU)}

\subsection{Motivation}

Many real-world data are sequential:

\begin{itemize}
    \item text (sentences, documents),
    \item speech signals,
    \item time series data.
\end{itemize}

\medskip

\noindent
\textbf{Key question.}  
How can we model data where order and temporal dependencies matter?

\medskip

\noindent
\textbf{Guiding idea.}  
Sequence models process inputs step-by-step, maintaining a hidden state that summarizes past information.

\subsection{Recurrent Neural Networks}

A recurrent neural network (RNN) updates a hidden state over time:

\[
h_t = \sigma(W_h h_{t-1} + W_x x_t + b),
\]

\[
y_t = W_y h_t.
\]

\medskip

\noindent
\textbf{Interpretation.}

\begin{itemize}
    \item $h_t$ stores information about the past
    \item The same parameters are reused across time steps
\end{itemize}

\medskip

\noindent
\textbf{Unrolled view.}

\begin{itemize}
    \item RNN can be viewed as a deep network with shared weights
    \item Depth corresponds to sequence length
\end{itemize}

\medskip

\noindent
\textbf{Insight.}  
RNNs extend neural networks to sequential data via parameter sharing over time.

\subsection{Long-Term Dependencies}

\textbf{Problem.}

\begin{itemize}
    \item Information from early time steps may be lost
    \item Gradients decay or explode over long sequences
\end{itemize}

\medskip

\noindent
\textbf{Mathematical intuition.}

\[
\frac{\partial L}{\partial h_t}
= \prod_{k=t}^{T} \frac{\partial h_{k+1}}{\partial h_k}.
\]

\medskip

\noindent
\textbf{Observation.}

\begin{itemize}
    \item Similar to deep networks in space, but now in time
    \item Suffers from vanishing/exploding gradients
\end{itemize}

\medskip

\noindent
\textbf{Insight.}  
Sequence length plays the role of depth in RNNs.

\subsection{Gated Architectures: LSTM}

Long Short-Term Memory (LSTM) networks introduce gates to control information flow.

\medskip

\noindent
\textbf{Core components.}

\begin{itemize}
    \item \textbf{Cell state} $c_t$: long-term memory
    \item \textbf{Gates:}
    \begin{itemize}
        \item input gate
        \item forget gate
        \item output gate
    \end{itemize}
\end{itemize}

\medskip

\noindent
\textbf{Update equations (conceptual).}

\begin{itemize}
    \item Forget:
    \[
    f_t = \sigma(W_f x_t + U_f h_{t-1})
    \]

    \item Input:
    \[
    i_t = \sigma(W_i x_t + U_i h_{t-1})
    \]

    \item Candidate:
    \[
    \tilde{c}_t = \tanh(W_c x_t + U_c h_{t-1})
    \]

    \item Cell update:
    \[
    c_t = f_t \odot c_{t-1} + i_t \odot \tilde{c}_t
    \]
\end{itemize}

\medskip

\noindent
\textbf{Key idea.}

\begin{itemize}
    \item gates regulate information flow
    \item prevent uncontrolled decay or explosion
\end{itemize}

\medskip

\noindent
\textbf{Insight.}  
LSTM introduces controlled memory dynamics.

\subsection{Gated Recurrent Units (GRU)}

GRU is a simplified gated architecture.

\medskip

\noindent
\textbf{Key components.}

\begin{itemize}
    \item update gate
    \item reset gate
\end{itemize}

\medskip

\noindent
\textbf{Characteristics.}

\begin{itemize}
    \item fewer parameters than LSTM
    \item similar performance in many tasks
\end{itemize}

\medskip

\noindent
\textbf{Insight.}  
GRU simplifies gating while retaining effectiveness.

\subsection{Limitations of Recurrence}

Despite improvements, RNN-based models have limitations.

\medskip

\noindent
\textbf{Sequential computation.}

\begin{itemize}
    \item cannot fully parallelize across time steps
\end{itemize}

\medskip

\noindent
\textbf{Long-range dependencies.}

\begin{itemize}
    \item still difficult to capture very long dependencies
\end{itemize}

\medskip

\noindent
\textbf{Training difficulty.}

\begin{itemize}
    \item optimization remains challenging
\end{itemize}

\medskip

\noindent
\textbf{Insight.}  
Recurrence introduces inherent computational and optimization constraints.

\subsection{Continuous-Time and Dynamical Systems Perspective (Advanced)}

Recurrent neural networks can be interpreted as discrete-time dynamical systems.

\medskip

\noindent
Recall the RNN update:
\[
h_t = \sigma(W_h h_{t-1} + W_x x_t + b).
\]

\medskip

\noindent
\textbf{Dynamical systems view.}

\begin{itemize}
    \item $h_t$ is the state of the system at time $t$
    \item the update defines a nonlinear map:
    \[
    h_t = F(h_{t-1}, x_t)
    \]
\end{itemize}

\medskip

\noindent
\textbf{Insight.}  
RNNs evolve a hidden state through time via repeated application of a nonlinear transformation.

\subsection{Connection to Discrete Dynamical Systems}

If inputs are fixed or absent, the system becomes autonomous:
\[
h_t = F(h_{t-1}).
\]

\medskip

\noindent
\textbf{Behavior.}

\begin{itemize}
    \item fixed points: $h^* = F(h^*)$
    \item stable vs unstable dynamics
    \item attractors in state space
\end{itemize}

\medskip

\noindent
\textbf{Interpretation.}  
RNNs can store information in stable states of the dynamics.

\subsection{Continuous-Time Limit}

We can interpret the recurrence as a discretization of a continuous-time system.

\medskip

\noindent
Consider:
\[
h_{t+1} = h_t + \Delta t \, F(h_t, x_t).
\]

\medskip

\noindent
In the limit $\Delta t \to 0$, this becomes an ordinary differential equation:
\[
\frac{dh(t)}{dt} = F(h(t), x(t)).
\]

\medskip

\noindent
\textbf{Interpretation.}

\begin{itemize}
    \item hidden state evolves continuously over time
    \item RNN is a discretized flow in state space
\end{itemize}

\medskip

\noindent
\textbf{Insight.}  
RNNs approximate continuous-time dynamical systems.

\subsection{Gradient Flow and Stability}

Backpropagation through time corresponds to sensitivity analysis of the dynamical system.

\medskip

\noindent
\textbf{Gradient evolution.}

\[
\frac{\partial L}{\partial h_t}
= \left( \prod_{k=t}^{T} \frac{\partial h_{k+1}}{\partial h_k} \right)
\frac{\partial L}{\partial h_T}.
\]

\medskip

\noindent
\textbf{Connection to dynamics.}

\begin{itemize}
    \item gradients depend on Jacobians of the system
    \item stability of dynamics determines gradient behavior
\end{itemize}

\medskip

\noindent
\textbf{Lyapunov intuition.}

\begin{itemize}
    \item contracting dynamics $\Rightarrow$ vanishing gradients
    \item expanding dynamics $\Rightarrow$ exploding gradients
\end{itemize}

\medskip

\noindent
\textbf{Insight.}  
Gradient instability reflects instability of the underlying dynamical system.

\subsection{Gating as Controlled Dynamics}

LSTM and GRU can be interpreted as modifying the dynamics.

\medskip

\noindent
\textbf{Example (LSTM cell).}

\[
c_t = f_t \odot c_{t-1} + i_t \odot \tilde{c}_t.
\]

\medskip

\noindent
\textbf{Interpretation.}

\begin{itemize}
    \item $f_t$ controls memory retention
    \item $i_t$ controls input injection
\end{itemize}

\medskip

\noindent
\textbf{Insight.}  
Gates act as adaptive coefficients controlling system stability.

\subsection{Connection to Residual Networks}

Recall residual networks:
\[
h_{k+1} = h_k + F(h_k).
\]

\medskip

\noindent
\textbf{Parallel.}

\begin{itemize}
    \item ResNet: depth as time (spatial layers)
    \item RNN: time as time (temporal steps)
\end{itemize}

\medskip

\noindent
\textbf{Unified view.}

\begin{itemize}
    \item both are discretizations of continuous flows
\end{itemize}

\medskip

\noindent
\textbf{Insight.}  
Deep networks and sequence models share a common dynamical systems foundation.

\subsection{Broader Connections}

This perspective connects to:

\begin{itemize}
    \item \textbf{Neural ODEs:} continuous-depth models
    \item \textbf{Control theory:} stability and feedback
    \item \textbf{Dynamical systems:} attractors and trajectories
\end{itemize}

\medskip

\noindent
\textbf{Key insight.}  
Learning in sequence models can be viewed as learning a dynamical system that evolves representations over time.

\subsection{A Unifying Perspective}

Sequence models extend neural networks by:

\begin{itemize}
    \item introducing temporal structure,
    \item sharing parameters across time,
    \item maintaining a hidden state.
\end{itemize}

\medskip

\noindent
\textbf{Core idea.}  
Learning involves propagating information through time as well as depth.

\medskip

\noindent
\textbf{Key insight.}  
RNNs treat time as another dimension of computation.

\subsection{Outlook}

In this lecture, we introduced:

\begin{itemize}
    \item recurrent neural networks,
    \item long-term dependency problems,
    \item gated architectures (LSTM, GRU),
    \item limitations of recurrence.
\end{itemize}

\medskip

\noindent
These limitations motivate alternative architectures.

\medskip

\noindent
In the next lecture, we study attention mechanisms and transformers, which overcome many limitations of recurrent models.

\medskip

\noindent
\textbf{Guiding principle.}  
Efficient sequence modeling requires balancing memory, computation, and parallelism.